% ==============================================================================
\newcommand{\ruta}[1]{\texttt{#1}}
% SECCIÓN 1: INTRODUCCIÓN Y JUSTIFICACIÓN ARQUITECTÓNICA
% ==============================================================================
\section{Diseño de Firmware}

El diseño del firmware para un sistema embebido robótico representa un desafío crítico: se deben ejecutar de manera concurrente tareas de control determinista en tiempo real junto con algoritmos pesados de procesamiento espacial y servicios no deterministas de comunicación inalámbrica.

Para evitar los cuellos de botella característicos de las arquitecturas tradicionales basadas en bucles secuenciales infinitos (\textit{Super-Loops}), se ha seleccionado una \textbf{Arquitectura Jerárquica en Capas} gobernada por el kernel de tiempo real \textbf{FreeRTOS} sobre la plataforma EDU-CIAA.

\subsection{Arquitectura propuesta y criterios de diseño}

El firmware se diseñará para integrar la percepción, la estimación de la pose,
la construcción del mapa y el control del rover diferencial. La arquitectura
será portable entre plataformas: el desarrollo y la validación funcional se
realizarán primero sobre ESP32 ya que no contamos con una EDU-CIAA en nuestras casas para desarrollar y probar software de aplicación; luego se portará a EDU-CIAA reutilizando la
aplicación y los drivers compartidos, y reemplazando la capa de adaptación al
hardware.

Se adopta una organización en tres capas para separar las operaciones propias
de cada placa de los algoritmos del robot. Esta separación permite probar los
módulos por partes, mantener una única implementación de la lógica del rover y
evitar que las decisiones de una placa condicionen el código común. La
arquitectura de ejecución se organizará en tareas periódicas controladas por
tiempo y planificadas por FreeRTOS, la implementación y caracteristicas de este todavia no se decidieron ya que primero optamos por validar los datos del hardware e integrarlo con una simple versión del scheduler.

El modelo elegido prioriza interfaces estables entre capas, unidades físicas
comunes y un main compartido. El período de muestreo y las prioridades se
ajustarán mediante mediciones de carga y latencia cuando se integren los
encoders, la construcción del mapa y el control de motores.

\newpage
\subsection{Las Tres Capas}

\begin{figure}[h!]
    \centering
    \includegraphics[width=1\textwidth]{imagenes/capas.png}
    \caption{Estructura jerárquica de la arquitectura en capas para el firmware de Arturito.}
    \label{fig:capas_software}
\end{figure}
\begin{description}

  \item[Capa 1 --- HAL (Hardware Abstraction Layer).]
       Será una interfaz común y tendra una implementación por placa. Su objetivo
        principal es abstraer el acceso a los periféricos y recursos de bajo
        nivel del microcontrolador, tales como GPIO, I2C, UART, PWM,
        temporizadores e interrupciones, evitando que las capas superiores
        dependan directamente de registros o características especificas de
        una arquitectura.

        Cada plataforma tendrá su propia implementación de estas funcionalidades.
        Por ejemplo, una implementación para la EDU-CIAA podrá configurar
        directamente registros del LPC4337, mientras que otra implementación
        podrá utilizar funciones provistas por sAPI o LPCOpen. De manera
        equivalente, una implementación para ESP32 podrá utilizar los drivers
        proporcionados por ESP-IDF.

        Por lo tanto, la existencia de la HAL no implica que todos los
        periféricos deban implementarse desde cero mediante acceso directo a
        registros. Si la arquitectura dispone de una biblioteca o driver que
        ya implementa dicha funcionalidad, la función correspondiente de la
        HAL podrá encapsular esa implementación. De esta manera, las capas
        superiores mantienen siempre la misma interfaz independientemente de
        cómo se resuelva internamente cada operación.

        La HAL no realizará tratamiento de las
        mediciones ni algoritmos propios del funcionamiento del rover.
        Sin embargo, en aquellos casos en que una biblioteca específica de
        una plataforma ya proporcione funciones de mayor nivel para acceder a
        un dispositivo, la HAL podrá encapsular dichas funciones con el fin de
        mantener una interfaz uniforme entre arquitecturas. En estos casos terminaria actuando como una HAL de perifericos.
        El procesamiento
        común de las mediciones continuará perteneciendo a los drivers de
        dispositivos. La HAL tampoco realizará integración de pose, mapeo ni
        navegación.

        \begin{visar}
Ejemplos de funciones HAL:
\begin{itemize}

  \item \texttt{hal\_gpio\_config()} 

  \item \texttt{hal\_gpio\_write()} 

  \item \texttt{hal\_i2c\_init()} 

  \item \texttt{hal\_i2c\_read()} 

  \item \texttt{hal\_pwm\_set\_duty()} 


  \item \texttt{hal\_imu\_get\_gyro\_z\_rads()} representa el caso en que la capa se comporta como HAL de perifericos. Si la plataforma dispone de una
        biblioteca específica para un periferico, podría encapsular directamente
        la función equivalente de dicha biblioteca. De esta forma, las capas
        superiores utilizan la misma interfaz independientemente de la
        implementación disponible en cada arquitectura.

\end{itemize}
        \end{visar}


  \item[Capa 2 --- Drivers de dispositivos.]
        Interpretará las señales físicas de cada dispositivo y entregará datos
        listos para la aplicación. Incluirá lectura, calibración, corrección y
        validación de datos de los sensores y podrá enviar comandos de
        generación de señales de control para el DRV8833. Sus interfaces no
        dependerán de Arduino ni de sAPI.

        \begin{visar}
Ejemplos de funciones de drivers:
\begin{itemize}

  \item \texttt{driver\_mpu6050\_get\_gyro\_z\_rads()}

  \item \texttt{driver\_vl53l0x\_get\_distance\_m()}

  \item \texttt{driver\_as5600\_get\_angle\_rad()} 

  \item \texttt{driver\_motor\_set\_pwm()}

  \item \texttt{driver\_motor\_set\_direction()}

\end{itemize}
        \end{visar}


  \item[Capa 3 --- Aplicación.]
        Contendrá la lógica propia del rover: adquisición y sincronización de
        muestras, cinemática diferencial y estimación de pose, actualización
        de la grilla de ocupación, selección de consignas según el modo y
        distribución de datos a la navegación y las comunicaciones. Acá
        vivirán los algoritmos de estimación de pose, mapeo y navegación.

        \begin{visar}
Ejemplos de funciones de aplicación:
\begin{itemize}

  \item \texttt{odometry\_update()} 



  \item \texttt{pose\_estimator\_update()}

  \item \texttt{occupancy\_grid\_update()} 

  \item \texttt{navigation\_update()} 

\end{itemize}
        \end{visar}


  \item[Infraestructura transversal --- FreeRTOS.]
        FreeRTOS no se considerará una capa adicional de la arquitectura, sino
        una infraestructura transversal utilizada para organizar la ejecución
        concurrente del sistema. El kernel se encargará de la planificación de
        las tareas, sus prioridades y temporización, además de proporcionar
        mecanismos de comunicación y sincronización.

        La lógica de aplicación podrá dividirse en tareas periódicas o
        activadas por eventos, por ejemplo para adquisición de sensores,
        control de velocidad, estimación de pose, actualización del mapa,
        navegación y comunicaciones. Cuando resulte necesario, los drivers y
        la HAL también podrán utilizar mecanismos proporcionados por FreeRTOS
        para sincronizar accesos a recursos compartidos.

        

\end{description}

\subsection{Software de PC e interfaz de usuario prevista}

El alcance para el programa de PC es visualizar la pose del rover y
el mapa actualizado. Para transmitir
pose y mapa sin cable se prevé incorporar Wi-Fi o Bluetooth. La elección entre ambos protocolos queda pendiente. Además desde allí se podran mandar órdenes START/STOP y, como objetivos
secundarios, teleoperación y planificación de rutas; su incorporación como
controles de esta interfaz queda por definir.

\begin{figure}[H]
    \centering
    \includegraphics[width=1\linewidth]{imagenes/InterfazDispExt.png}
    \caption{Interfaz para ver el mapeo desde el punto de vista del usuario.}
    \label{fig:interfaz-dispext}
\end{figure}

Desde esta vista se puede imaginar el siguiente ciclo: 

\begin{center}
\fbox{\parbox{0.91\linewidth}{\centering
Inicio $\rightarrow$ conexión con el rover $\rightarrow$ espera de datos
$\rightarrow$ visualización continua de pose y mapa $\rightarrow$ pérdida de
enlace: indicar desconexión / conexión restablecida: reanudar visualización}}
\end{center}

Si se pierde el enlace durante la operación, el firmware deberá aplicar la
parada segura y la interfaz indicará la desconexión. 

\subsection{Bibliotecas y herramientas externas previstas}

\begin{itemize}
  \item \textbf{ESP32:} framework Arduino y su integración con FreeRTOS;
        \texttt{Wire} para I2C; biblioteca Adafruit \texttt{VL53L0X} para el
        sensor ToF. Para el transporte inalámbrico se utilizará la capacidad
        disponible en ESP32 una vez elegido el protocolo de enlace.
        Para la compilación, carga y estructura general del sistema se utiliza PlatformIO.
  \item \textbf{EDU-CIAA:} LPCOpen y SAPI para la placa y sus periféricos, más
        el kernel FreeRTOS configurado en el proyecto. La HAL de algunos modulos estan implementados (por ejemplo: MPU60x0, el cual brinda acceso simplificado al sensor MPU6050) mientras que hay sensores como el ToF que no tienen esta libreria especifica y al cual le tendremos que desarrollar la implementación de acceso I2C por registros.

  \item \textbf{Workspace:} Se logró integrar ambos entornos de desarrollo en un mismo wokspace para VScode. 

  \begin{figure}[h!]
    \centering
    \includegraphics[width=0.3\textwidth]{imagenes/workspace.png}
    \caption{Estrucura del Workspace.}
    \label{fig:workspace}
\end{figure}
  
  \begin{description}
  \item[\ruta{docs/}] Documentos del equipo.
  \item[\ruta{shared/}] Contratos HAL, drivers y logica de aplicación que no
        deben depender de Arduino ni de sAPI.
  \item[\ruta{esp32/}] Firmware PlatformIO para ESP32, sus HAL de placa y
        programas de diagnostico independientes.
  \item[\ruta{edu-ciaa/}] Copia de estructura del entorno de 'Embedded IDE'. Make, LPCOpen, sAPI, todo lo necesario para la compilación y carga sobre EDU-CIAA. Adapta la
        misma aplicación y drivers a su arquitectura.
  \item[\ruta{processing/}] Por ahora esta carpeta es nuestra "PC externa". Contiene sketches que reciben datos para visualizar.
\end{description}
\end{itemize}


\newpage

\subsection{Versiones y ensayos realizado}
\subsubsection{Pruebas IMU y ToF}
 Esta versión ya cuenta con la estructura base del workspace explicado anteriormente pero solo realiza el siguiente ciclo de percepción:

\begin{enumerate}
  \item Inicializa la IMU MPU6050 y el sensor de distancia VL53L0X.
  \item Lee la velocidad angular del eje Z de la IMU y una distancia del ToF.
  \item Corrige e integra la velocidad angular para estimar el angulo $\theta$.
  \item Acepta distancias entre 0.05 m y 1.50 m; el resto se representa como
        $-1$ (medida invalida o sensor no disponible).
  \item Publica una linea CSV con $\theta$ en radianes y distancia en metros.
  \item Processing (programa graficador) representa la orientación y acumula puntos del ToF mientras
        el robot gira en el lugar.
\end{enumerate}

No hay todavía odometria XY, encoders, motores, navegación autónoma ni grilla de
ocupación. No hay modos de funcionamiento por comandados ni comunicación inalámbrica.

El firmware fue cargado tanto en ESP32 como en EDU-CIAA y validado con Processing: se comprobó la
recepción de las muestras serie y la visualización de orientación y distancia.
\begin{figure}[h!]
    \centering
    \includegraphics[width=0.6\textwidth]{imagenes/V0.png}
    \hfill
    \includegraphics[width=0.38\textwidth]{imagenes/f1.png}
    \caption{Visualización en Processing (izquierda) y sensores montados en un auto a control remoto (derecha). \href{https://drive.google.com/file/d/1k8-xx6FHV0KANsshld4rqhePZVvpV37A/view?usp=drive_link}{Link para ver funcionamiento en video.}}
    \label{fig:v0_y_sensores}
\end{figure}

\newpage

Cómo llegamos a este punto:

El Firmware es en realidad una versión final de un proceso de testeo de los sensores en ambas arquitecturas, de un proceso de interpretación de datos, de armado de workspace, de ejecución del programa con FreeRTOS, etc.

Algunos cosas interesantes de mensionar:
\begin{description}
 \item[Angulo entregado por capa 1 vs Angulo provisto por capa 2:]
 
 La capa 2 espera que el robot quede quieto, promedia 1000 lecturas separadas por 2
ms para estimar el offset, y reinicia el estado de integración, luego resta el offset, elimina velocidades dentro de ±1◦/s, calcula un promedio trapezoidal entre muestras y acumula el angulo con el tiempo transcurrido.
   \begin{figure}[h!]
    \centering
    \includegraphics[width=0.9\textwidth]{imagenes/V0-grados.jpeg}
    \caption{Obtención del angulo crudo vs. procesado por capa 2.}
    \label{fig:grados_crudos}
\end{figure}

Se ve que a tan solo 30s de inicializado el rover el angulo crudo difiere en 5° del angulo real (el rover estaba quieto).

   \begin{figure}[h!]
    \centering
    \includegraphics[width=0.9\textwidth]{imagenes/resta.jpeg}
    \caption{Angulo procesado - Angulo crudo}
    \label{fig:resta_grados}
\end{figure}
Acá se puede ver que si se mueve el robot cambia la relacion entre una señal y la otra. Si el procesado no fuese necesario entonces este grafico deberia ser casi siempre 0 o aunque sea aproximarse a 0. 

 \item[Error en la librería de MPU para EDU-CIAA:] como ya mencionamos, existe una librería que otorga recursos para manejar el giroscopio y acelerometro MPU6050 dentro de sAPI. Durante las pruebas de integración en hardware de se identificó un error que hacia que el sensor no responda.

El procedimiento de arranque de la libreria implementaba una verificación de integridad mediante la lectura del registro interno WHO\_AM\_I (dirección 0x75), cuyo valor de retorno esperado según la hoja de datos estándar del fabricante es 0x68. Sin embargo, la función de inicialización abortaba de forma consistente retornando un código de falla (-5), impidiendo la ejecución del resto del sistema.

Tras analizar el bus de comunicación I2C mediante rutinas de escaneo, se constató que la variante física del módulo utilizado responde con el identificador 0x72 no 0x68. Entendimos que este comportamiento es habitual en hardware genéricos y para solucionar esta incompatibilidad se reestructuró la validación en el driver. La lógica fue modificada para admitir de forma flexible tanto el identificador nativo (0x68) como el valor identificado en el banco de pruebas (0x72).

\end{description}


\subsubsection{Estructura de compilación y ejecución de aplicaciones}

A medida que avanzó el desarrollo del firmware, se incorporaron distintas aplicaciones de prueba para verificar individualmente los sensores, actuadores y módulos del sistema, antes de integrarlos en el funcionamiento general del rover.

Uno de los objetivos fue mantener la compatibilidad con ambas arquitecturas, ESP32 y EDU-CIAA, evitando desarrollar dos versiones independientes de una misma aplicación. Para ello, se implementó un mecanismo común de inicialización y ejecución, adaptado a los requisitos particulares de cada plataforma.

En el caso del ESP32, utilizando el framework Arduino, la ejecución requiere las funciones \texttt{setup()} y \texttt{loop()}, destinadas a la inicialización y ejecución repetitiva del programa, respectivamente. Por otro lado, la EDU-CIAA utiliza una función \texttt{main()}, desde la cual se inicializan los recursos de la placa y se controla el ciclo principal de ejecución.

Para unificar ambos mecanismos se definieron dos funciones comunes:

\begin{itemize}
	\item \texttt{app\_init():} realiza la inicialización de los recursos y módulos necesarios para la aplicación seleccionada.
	\item \texttt{app\_update():} ejecuta una iteración del funcionamiento de dicha aplicación.
\end{itemize}

Cada arquitectura dispone de un archivo de entrada propio (\texttt{entry\_esp32.cpp} y \texttt{entry\_educiaa.c}), encargado de inicializar los recursos generales de la plataforma y llamar a estas dos funciones. En ESP32, las llamadas se realizan desde \texttt{setup()} y \texttt{loop()}, mientras que en EDU-CIAA se efectúan desde \texttt{main()}, utilizando un ciclo \texttt{while(1)}.

De esta manera, las aplicaciones de capa 3 no necesitan conocer el mecanismo de arranque de la placa donde se ejecutan. Las diferencias de hardware continúan resolviéndose mediante las implementaciones específicas de la HAL, conservando la separación por capas definida anteriormente.

\paragraph{Selección y compilación de aplicaciones:}

Cada prueba se implementa como una aplicación independiente que define las funciones \texttt{app\_init()} y \texttt{app\_update()}. Para ejecutarla, se seleccionan durante la compilación los archivos correspondientes a dicha aplicación, junto con los drivers y las implementaciones de HAL que necesita.

El procedimiento de selección depende del entorno de desarrollo:

\begin{description}
	\item[ESP32:] Se utiliza PlatformIO, donde el archivo \texttt{platformio.ini} contiene distintos entornos de compilación. Cada entorno especifica los archivos fuente y las dependencias necesarias mediante \texttt{build\_src\_filter}, además de heredar una configuración común de la placa. Para ejecutar una prueba diferente, basta con seleccionar su entorno correspondiente.
	
	\item[EDU-CIAA:] Se utiliza el sistema de compilación basado en Make. El archivo \texttt{program.mk} define los modos de aplicación disponibles mediante la variable \texttt{APP\_MODE}, mientras que \texttt{rover/config.mk} selecciona los archivos fuente y las dependencias correspondientes. El modo puede indicarse directamente al ejecutar el comando \texttt{make}, sin necesidad de modificar el \texttt{Makefile} principal.
\end{description}

Esta organización permite realizar pruebas individuales de la IMU, el sensor ToF, los encoders, los motores y otros módulos, así como desarrollar ensayos que integren varios componentes.

La principal ventaja es que una misma aplicación de prueba puede reutilizarse en ambas arquitecturas, cambiando únicamente el entorno de compilación y las implementaciones de hardware utilizadas. Además, permite incorporar nuevas pruebas sin modificar los puntos de entrada ni alterar las aplicaciones existentes, facilitando el desarrollo incremental y la detección de errores durante la integración.

\subsubsection{Sistema de telemetría}

Para facilitar las pruebas y el diagnóstico de los distintos módulos, se desarrolló un sistema de telemetría que permite transmitir información desde el rover hacia una computadora externa.

Se buscó que las aplicaciones pudieran enviar información independientemente de la arquitectura utilizada, manteniendo una interfaz común y delegando las particularidades de comunicación a las implementaciones de HAL de cada placa.

El sistema utiliza el ESP32 como encargado de establecer la comunicación inalámbrica mediante Wi-Fi. Actualmente funciona en modo Access Point (AP), creando una red propia denominada \texttt{Arturito}, a la cual puede conectarse una computadora sin necesidad de contar con un router externo. Sobre esta conexión se utiliza un servidor TCP, que transmite los mensajes de telemetría a través del puerto 3333.

La implementación presenta dos modalidades según la arquitectura que ejecuta la aplicación:

\begin{description}
	
	\item[Aplicación ejecutada en ESP32:]
	La aplicación utiliza la interfaz de telemetría compartida, cuya implementación de HAL envía los mensajes directamente a través del módulo Wi-Fi del ESP32. De esta manera, un único firmware realiza tanto el procesamiento de la aplicación como la transmisión de sus resultados.
	
	\item[Aplicación ejecutada en EDU-CIAA:]
	En este caso, la EDU-CIAA no dispone de conectividad Wi-Fi propia, por lo que utiliza una comunicación UART para transmitir los mensajes hacia el ESP32.
	
	Para ello, se carga en el ESP32 un firmware independiente que funciona como puente UART--Wi-Fi. Este programa recibe los mensajes enviados por EDU-CIAA y los retransmite mediante la conexión TCP, sin realizar el procesamiento de los datos recibidos.
	
\end{description}

En ambos casos, las aplicaciones utilizan las mismas funciones de telemetría, sin necesidad de conocer el medio físico de transmisión. Esto permite mantener compartido el código de las pruebas y modificar la arquitectura utilizada sin alterar su funcionamiento general.

La transmisión se realiza mediante mensajes de texto delimitados por saltos de línea, lo que facilita su recepción e interpretación desde una computadora.

Se desarrolló una prueba específica de telemetría que transmite mensajes periódicos, con el objetivo de verificar la comunicación y el funcionamiento de ambas modalidades.

   \begin{figure}[h!]
	\centering
	\includegraphics[width=0.9\textwidth]{imagenes/wifi.png}
	\caption{Telemetría para ambas arquitecturas}
	\label{fig:wifi}
\end{figure}

\subsubsection{Pruebas de sensores y actuadores}

Una vez establecida la estructura común de ejecución de aplicaciones y el sistema de telemetría, se desarrollaron distintas pruebas para verificar el funcionamiento de los sensores y actuadores del rover.

Estas pruebas permitieron evaluar los módulos de forma independiente antes de avanzar hacia su integración. En todos los casos se mantuvo la arquitectura por capas, utilizando los mismos drivers y lógica de aplicación para ambas plataformas, con las correspondientes implementaciones de HAL.

\begin{description}
	
	\item[Prueba de IMU (MPU6050):]
	Se desarrolló una aplicación que inicializa y calibra el sensor, solicitando que el rover permanezca quieto durante este procedimiento.
	
	Una vez finalizada la calibración, se realizan lecturas periódicas de la velocidad angular para estimar el ángulo de orientación alrededor del eje Z. El ángulo calculado es transmitido mediante telemetría, permitiendo observar su comportamiento mientras se gira físicamente el sensor.
	
	Esta prueba permite verificar tanto la comunicación con la IMU como el funcionamiento del procesamiento de orientación implementado en las capas superiores.
	
	\item[Prueba de distancia (VL53L0X):]
	Se implementó una aplicación para verificar las mediciones del sensor ToF mediante comunicación I2C.
	
	El sensor se inicializa utilizando el perfil de medición High Speed y se realizan lecturas periódicas de distancia. Para cada medición se transmiten dos valores: la distancia obtenida directamente del sensor y la distancia corregida mediante la función de calibración implementada.
	
	De esta manera, es posible comparar ambas mediciones y evaluar el comportamiento de la corrección aplicada. La prueba también contempla la detección de mediciones inválidas y errores de comunicación.
	
	\item[Pruebas de encoders (AS5600):]
	Se desarrollaron pruebas para verificar la comunicación I2C con los encoders magnéticos instalados en las ruedas del rover.
	
	Inicialmente se realizaron lecturas de la posición angular absoluta de cada encoder, girando manualmente las ruedas y transmitiendo los valores obtenidos.
	
	Posteriormente, estas lecturas se utilizaron para calcular la velocidad angular a partir de la variación de posición y el tiempo transcurrido entre muestras, contemplando el salto de ángulo que se produce al completar una revolución.
	
	\item[Pruebas de motores (DRV8833):]
	Se implementaron aplicaciones para comprobar el accionamiento de los motores mediante señales PWM generadas a través de la HAL correspondiente.
	
	Se ensayaron diferentes secuencias de movimiento, incluyendo avance, retroceso y giros sobre el propio eje. También se desarrollaron pruebas de trayectorias predefinidas, utilizando tiempos de actuación y diferentes valores de PWM.
	
	Estas primeras pruebas se realizaron en lazo abierto, sin utilizar realimentación de los encoders para corregir el movimiento.
	
	\item[Integración de motores y encoders:]
	Finalmente, se combinaron las pruebas anteriores para accionar ambos motores y obtener simultáneamente información de los encoders.
	
	Se implementaron secuencias con distintos sentidos de giro y valores de PWM, mientras se calculaba y transmitía la velocidad angular de las ruedas en RPM.
	
	Esta integración permite relacionar la señal de control aplicada a los motores con la velocidad medida, constituyendo una primera base experimental para el desarrollo posterior del control de velocidad.
	
\end{description}

Los resultados y mensajes de diagnóstico de estas pruebas se transmitieron mediante el sistema de telemetría, permitiendo observar el comportamiento de los distintos módulos desde una computadora externa.
